\chapter{Trabalhos Relacionados}
\label{chap:relacionados}

\epigraph{\itshape ``Se eu vi mais longe, foi por estar sobre ombros de gigantes.''}
{\normalfont --- Isaac Newton, carta a Robert Hooke, 1675}

A revisão da literatura organiza-se em sete eixos que convergem para o problema
central desta dissertação: (i) frameworks de navegação autônoma em ROS~2, (ii)
avaliação e comparação de algoritmos de SLAM, (iii) metodologias de avaliação
de repetibilidade em robótica, (iv) sistemas de orquestração para frotas de
robôs, (v) interfaces web para operação de robôs, (vi) Engenharia de Software
aplicada a sistemas robóticos baseados em ROS, e (vii) a crise de reprodutibilidade
em pesquisa empírica de Engenharia de Software e o papel dos artefatos de pesquisa.
A análise de cada eixo identifica contribuições relevantes e lacunas que esta
dissertação pretende preencher.

\section{Frameworks de Navegação Autônoma em ROS~2}
\label{sec:rel_nav}

\subsection{Navigation2 (Nav2) e o Marathon 2}

O trabalho mais diretamente relacionado a esta dissertação é a introdução do
Nav2 no contexto do \textit{Marathon 2} \cite{macenski2020marathon}. O artigo
descreve a arquitetura do Nav2 --- servidores de planejamento, controle e
comportamento orquestrados por Behavior Trees --- e valida a pilha em mais de
60\,km (37,4 milhas) de navegação autônoma acumulada em um campus universitário
dinâmico e povoado por pessoas, usando dois robôs profissionais (Tiago e RB-1),
sem nenhuma colisão registrada ao longo de quase 23 horas de operação.

A contribuição é fundamental: demonstra que o Nav2 é maduro o suficiente para
uso contínuo em ambientes não estruturados. No entanto, o critério de avaliação
é agregado (distância total percorrida, número de intervenções humanas) e não
inclui análise de consistência entre execuções individuais. O framework proposto
nesta dissertação não questiona a qualidade do Nav2 --- constrói sobre ela para
responder perguntas que o Marathon 2 não se propôs a responder.

\subsection{O Controlador Regulated Pure Pursuit}

\citeonline{macenski2023survey} apresentam o controlador Regulated Pure Pursuit para
ROS~2, uma adaptação do clássico Pure Pursuit \cite{coulter1992} com regulação
de velocidade adaptativa. O artigo demonstra rastreamento de trajetória com
erro lateral médio abaixo de 5\,cm em curvas acentuadas --- uma métrica diretamente
relacionada à repetibilidade de percurso.

A diferença em relação a esta dissertação é de escopo: o artigo avalia o desempenho
de rastreamento de um controlador específico, enquanto esta dissertação avalia a
repetibilidade do sistema completo (SLAM + Nav2 + odometria) como um pipeline
integrado. O erro de rastreamento do controlador é apenas uma das contribuições
para o RMSE total observado.

\subsection{O Ecossistema ROS~2 em Aplicações Reais}

\citeonline{macenski2022ros2} realizam uma análise abrangente do ROS~2 cerca de
cinco anos após seu primeiro lançamento estável (Ardent Apalone, dezembro de 2017).
O artigo documenta cinco estudos de caso reais --- um robô terrestre com pernas
(Ghost Robotics), uma embarcação de superfície autônoma para missões marinhas
(Mission Robotics), um sistema de drones aéreos (Auterion), o rover lunar VIPER
da NASA, e a orquestração de frotas industriais em larga escala (OTTO Motors)
--- demonstrando que o middleware escalou de um framework experimental para
infraestrutura de produção em domínios críticos e heterogêneos. Para esta
dissertação, o artigo é a referência primária para o modelo de QoS, o protocolo
de lifecycle e a arquitetura do DDS.

\section{Avaliação de Algoritmos de SLAM}
\label{sec:rel_slam}

\subsection{A Taxonomia do SLAM}

O tutorial em duas partes de \citeonline{durrantwhyte2006} e \citeonline{bailey2006}
consolidou e popularizou a taxonomia do problema SLAM --- cuja estrutura, resultado
de convergência e nomenclatura haviam sido propostos originalmente por Durrant-Whyte,
Rye e Nebot em 1996 --- e forneceu uma análise sistemática abrangente das abordagens
então existentes. Vinte anos depois, essa taxonomia permanece válida: filtros
de partículas (Gmapping, AMCL), filtros de Kalman estendidos (EKF-SLAM), e grafos
de poses (SLAM Toolbox, Cartographer) ainda são as três famílias algorítmicas
dominantes \cite{thrun2005}.

\subsection{Avaliação Empírica de SLAM em ROS}

\citeonline{hamer2025} realizam uma avaliação empírica rigorosa de quatro algoritmos
de SLAM 2D (Cartographer, Gmapping, Hector SLAM \cite{kohlbrecher2011} e Karto) em
sistemas baseados em ROS, medindo consumo de energia, utilização de CPU e acurácia de mapa ao longo de
múltiplas execuções em que o mesmo robô alterna entre algoritmos em cenários
ponto-a-ponto e circulares. O ponto metodológico mais relevante para esta
dissertação é a utilização de múltiplas execuções do mesmo cenário, com análise
de variabilidade entre execuções --- um embrião de análise de repetibilidade,
embora sem um protocolo formal de record--replay nem métricas de RMSE entre
trajetórias.

O estudo é publicado no \textit{Journal of the Brazilian Computer Society}, reforçando
esse periódico como veículo relevante para avaliações empíricas rigorosas de sistemas
baseados em ROS e a comparação empírica como metodologia científica consolidada na área.

\subsection{Repetibilidade em Mapeamento 3D}

\citeonline{maset2022} avaliam a repetibilidade de um sistema robótico de mapeamento
3D (scanner a laser Velodyne acoplado a uma unidade de medição inercial sobre uma
plataforma móvel), comparando três aquisições robóticas repetidas do mesmo ambiente
entre si, com três aquisições manuais (\textit{handheld}) e com um levantamento a
laser terrestre usado como referência. As métricas principais são a distância
\textit{cloud-to-cloud} entre nuvens de pontos e o RMSE posicional; a distância
média entre nuvens da mesma modalidade fica na casa de 2\,cm, com o modo robótico
apresentando ruído ligeiramente menor que o modo manual.

Embora focado em mapeamento 3D com hardware de custo elevado, o artigo estabelece
um conceito metodológico que esta dissertação adapta para navegação 2D: avaliar
repetibilidade por comparação direta entre execuções nominalmente idênticas do
mesmo percurso, em vez de comparar cada execução apenas contra um valor de
referência único --- e a importância de reportar não apenas a média dos erros,
mas também o desvio-padrão entre execuções.

\section{Avaliação de Repetibilidade em Robótica}
\label{sec:rel_repeticao}

\subsection{Boas Práticas Metodológicas}

\citeonline{amigoni2014} realizam um levantamento sistemático de como experimentos são
conduzidos em artigos de robótica autônoma publicados na \textit{International Conference
on Autonomous Agents and Multiagent Systems} (AAMAS) entre 2002 e 2012. Os resultados
mostram que boa parte dos trabalhos ainda carece de rigor científico --- comparação
estatística fraca (presente em pouco mais de um terço dos artigos que reportam apenas
médias), baixa disponibilidade de código e dados publicados (a maioria dos links
fornecidos está quebrada) e uso predominante de comparações apenas com variantes do
próprio sistema proposto, não com sistemas alternativos. Esse diagnóstico, publicado
em 2014, permanece pertinente em 2026: a ausência de protocolos padronizados de
avaliação continua sendo um obstáculo para o progresso cumulativo da área.

\citeonline{amigoni2010} também propõem que simulação e experimentação física devem ser
tratadas como complementares, não excludentes: a simulação permite controle completo
das variáveis, enquanto o hardware real captura fenômenos físicos irreplicáveis em
simulação. Esta dissertação adota essa visão ao focar em simulação para validação
do framework, identificando hardware real como extensão natural.

\subsection{Fundamentos Probabilísticos da Avaliação de Navegação}

\citeonline{thrun2005}, em seu tratado sobre robótica probabilística, não dedicam um
capítulo isolado à avaliação de navegação, mas tratam a incerteza de localização e a
convergência de estimativas como tema central ao longo dos capítulos sobre filtros de
Bayes, filtros de Kalman, filtros de partículas e SLAM (incluindo o FastSLAM). Essa
formalização recorrente do erro de localização esperado e da convergência de
estimativas fornece o embasamento teórico para as métricas de RMSE utilizadas nesta
dissertação. O livro, embora publicado em 2005, permanece a referência mais citada em
robótica probabilística e sua relevância para avaliar navegação com SLAM é
indiscutível.

\section{Sistemas de Orquestração para Frotas}
\label{sec:rel_frota}

\subsection{Arquiteturas de Frota em ROS~2}

O ROS~2 provê os primitivos de comunicação necessários para sistemas multi-robô,
mas não uma camada de orquestração de alto nível \cite{macenski2022ros2}. Propostas
como o Fleet Management do Open RMF (Open Robotics Middleware Framework) abordam o
problema de despacho de tarefas em frotas com múltiplos robôs, mas com foco em
planejamento de missões e otimização de rotas, não em repetibilidade experimental.

A distinção é fundamental e define o posicionamento desta dissertação: um sistema
de gestão de frota maximiza a eficiência operacional; o framework proposto aqui
garante que cada robô execute exatamente o que foi solicitado, nas mesmas condições,
toda vez. Os dois objetivos são ortogonais e complementares.

\subsection{Comportamento Determinístico em Sistemas Multi-Robô}

A dificuldade de garantir comportamento determinístico em sistemas multi-robô
com ROS~2 é documentada em trabalhos de pesquisa de tempo-real. A não-determinismo
no scheduling do sistema operacional, o jitter do DDS e as variações de tempo de
processamento entre execuções contribuem para a variabilidade observada em
experimentos de navegação \cite{macenski2022ros2}.

O framework proposto mitiga esse problema com estratégias de software --- espera
explícita pela ativação do Nav2, cooldown entre fases do experimento, e filtro de
saltos impossíveis na odometria --- mas não elimina a fonte fundamental do
não-determinismo, que reside no sistema operacional de tempo compartilhado.

\section{Interfaces Web para Robótica}
\label{sec:rel_web}

A maioria das interfaces de operação para robôs autônomos é baseada no RViz2,
uma ferramenta de visualização que requer instalação local do ROS~2 e não oferece
acesso remoto nativo. Alternativas web existem --- Foxglove Studio, Robot Web
Tools \cite{toris2015} --- mas requerem o ROSBridge WebSocket server, adicionando
uma camada de configuração e um ponto de falha ao sistema. O próprio Robot Web
Tools, na origem do protocolo rosbridge amplamente adotado por essas alternativas,
foi projetado como uma arquitetura cliente-servidor voltada à robótica em nuvem e
à interoperabilidade entre navegadores e múltiplos robôs \cite{toris2015}, um
objetivo mais amplo do que a interface de operação local e single-robot proposta
nesta dissertação.

A abordagem adotada nesta dissertação --- FastAPI como bridge REST sem ROSBridge,
com canvas HTML5 para visualização do mapa --- representa uma solução de menor
peso de configuração. A ausência de trabalhos publicados que adotem especificamente
essa combinação (FastAPI + canvas + Nav2 sem ROSBridge) reforça a originalidade
da contribuição de interface.

\section{Engenharia de Software em Sistemas Robóticos Baseados em ROS}
\label{sec:rel_es}

Embora os eixos anteriores situem esta dissertação no espaço técnico da robótica
autônoma, sua motivação central --- infraestrutura reprodutível, verificável e
auditável para experimentos --- é, antes de tudo, um problema de Engenharia de
Software aplicado a um domínio físico. Esta seção situa o trabalho nessa literatura.

\citeonline{albonico2023} conduzem um mapeamento sistemático sobre pesquisa em
Engenharia de Software voltada ao ROS, cobrindo 63 estudos primários. Os autores
identificam que a relevância industrial das avaliações propostas na literatura é
predominantemente baixa (47 dos 63 estudos) e que o rigor metodológico falha
sobretudo no relato de ameaças à validade (46 dos 63 estudos não discutem esse
aspecto), enquanto 19 estudos reconhecem explicitamente a necessidade de validação
adicional e 15 relatam falta de generalização de seus próprios resultados. Os
autores apontam ainda que o ROS~2 permanece pouco explorado tanto na pesquisa
(apenas 8 dos 63 estudos) quanto na adoção industrial, que segue limitada apesar
do fim do suporte de longo prazo ao ROS~1 previsto para 2025. Esse mapeamento
fundamenta diretamente o problema central desta dissertação: a crise de
orquestração experimental (Seção~\ref{sec:problema}) é, na taxonomia de
\citeauthor{albonico2023}, um sintoma da ausência de infraestrutura de teste,
validação e reprodutibilidade madura para sistemas ROS.

\citeonline{timperley2024} apresentam o ROBUST, um levantamento empírico de 221 bugs
reais coletados de sete sistemas de software baseados em ROS~1 (kobuki, mavros,
geometry, turtlebot, motoman, entre outros), catalogando causas-raiz que incluem
falhas de concorrência (ausência ou uso incorreto de primitivas de sincronização
entre nós), configuração incorreta de comunicação em rede (tamanhos de fila de
publicador mal configurados, causando bloqueios não intencionais) e uma classe
ampla de erros de programação geral. Embora o dataset seja anterior ao ROS~2 e,
portanto, não trate diretamente de conceitos específicos dessa geração do
middleware --- como QoS ou o modelo de \textit{executor} --- as categorias de
falha de concorrência e má configuração de comunicação identificadas por
\citeauthor{timperley2024} são análogas, em espírito, a problemas de
compatibilidade enfrentados durante o desenvolvimento do framework proposto
(Seção~\ref{sec:impl_orquestrador}), como o mismatch de QoS entre o coletor de
dados e o SLAM Toolbox e os deadlocks do executor de serviço único frente ao
\texttt{wait\_\allowbreak for\_\allowbreak server} do Nav2 --- manifestações específicas do ROS~2 da
mesma classe geral de defeito de concorrência/comunicação já documentada no
ecossistema ROS~1. Isso reforça que tais problemas não são peculiaridades de
uma implementação isolada, mas uma classe de defeito recorrente e não-trivial
que atravessa gerações do middleware.

\citeonline{santos2024} discutem a integração de práticas de monitoramento de
Integração Contínua (CI) em pipelines de engenharia de software, com o Holter
como ferramenta de observabilidade de builds. Embora o domínio de aplicação seja
distinto (CI de software convencional, não robótica), o princípio de tornar
comportamentos regressivos observáveis automaticamente --- em vez de depender de
inspeção manual --- é o mesmo princípio que motiva o protocolo record--replay com
análise RMSE automatizada proposto nesta dissertação: cada execução do protocolo
gera um oráculo quantitativo (o relatório de repetibilidade) que pode, em
princípio, ser incorporado a um pipeline de CI/CD para navegação autônoma, tal
como sugerido na agenda de trabalhos futuros (Seção~\ref{sec:futuros}).

Em conjunto, esses três trabalhos deslocam o enquadramento desta dissertação: o
framework proposto não deve ser lido apenas como uma ferramenta de robótica que
incidentalmente automatiza tarefas repetitivas, mas como uma resposta direta a uma
lacuna documentada na literatura de Engenharia de Software sobre sistemas ROS ---
a ausência de infraestrutura reutilizável de teste, observabilidade e
reprodutibilidade experimental.


\section{A Crise de Reprodutibilidade e o Papel dos Artefatos de Pesquisa}
\label{sec:rel_reprodutibilidade}

Um último eixo, transversal aos anteriores, situa esta dissertação no debate mais
amplo sobre reprodutibilidade em pesquisa empírica de Engenharia de Software.
\citeonline{guevaravega2024} analisam 106 estudos primários envolvendo experimentos
orientados a humanos e propõem uma estrutura de artefato de pesquisa alinhada à
política de \textit{badges} da ACM (\textit{Artifacts Available},
\textit{Artifacts Evaluated -- Functional} e \textit{Results Replicated}), com o
objetivo de tornar experimentos publicados de fato reexecutáveis por terceiros ---
não apenas descritos em texto. Os autores relatam que uma fração substancial dos
artefatos de pesquisa disponibilizados por autores de estudos empíricos está
incompleta ou imprecisa o suficiente para impedir a reprodução direta dos
resultados, mesmo quando código e dados estão nominalmente publicados.

O framework proposto nesta dissertação pode ser lido, sob essa lente, como uma
tentativa de aplicar o princípio dos \textit{badges} ACM \emph{durante a coleta
dos dados}, e não apenas na publicação posterior do artefato: cada execução do
protocolo record--replay (Capítulo~\ref{chap:metodologia}) produz, automaticamente,
um \texttt{summary.json} com as métricas RMSE pairwise, os CSVs de trajetória por
execução e o gráfico de sobreposição (\texttt{trajectory\_\allowbreak overlay.png}) descritos
na Seção~\ref{sec:impl_pipeline} --- artefatos estruturados o suficiente para que
outro pesquisador, de posse do mesmo \textit{workspace} ROS~2 e do mesmo mundo de
simulação, consiga reexecutar o experimento e comparar diretamente os resultados,
sem depender de reconstrução manual do protocolo a partir da prosa do artigo.
Essa característica é precisamente o tipo de lacuna que motivou a rejeição do
artigo original submetido à trilha de ferramentas da SBES (Apêndice~\ref{ap:fleet_msgs}
documenta a interface que viabiliza essa automação): a contribuição de
Engenharia de Software do framework não está nos algoritmos de navegação em si
--- que são reaproveitados do Nav2 e do SLAM Toolbox sem modificação --- mas na
infraestrutura que torna cada execução um artefato de pesquisa autocontido.

\section{Síntese Comparativa}
\label{sec:rel_sintese}

A Tabela~\ref{tab:comparacao_relacionados} resume os trabalhos discutidos e
compara suas características com o framework proposto nesta dissertação.

\begin{table}[H]
\centering
\caption{Comparação do framework proposto com trabalhos relacionados}
\label{tab:comparacao_relacionados}
\footnotesize
\begin{tabularx}{\textwidth}{lXccccc}
\toprule
\textbf{Trabalho} & \textbf{Foco principal} & \rotatebox{90}{\textbf{Record-Replay}} & \rotatebox{90}{\textbf{Multi-robô}} & \rotatebox{90}{\textbf{Web UI}} & \rotatebox{90}{\textbf{Métricas RMSE}} & \rotatebox{90}{\textbf{Open-source}} \\
\midrule
\citeonline{macenski2020marathon}  & Navegação longa distância & $\times$ & $\times$ & $\times$ & $\times$ & \checkmark \\
\citeonline{macenski2023survey}    & Controlador Pure Pursuit  & $\times$ & $\times$ & $\times$ & Parcial  & \checkmark \\
\citeonline{macenski2022ros2}      & Plataforma ROS~2          & $\times$ & Parcial  & $\times$ & $\times$ & \checkmark \\
\citeonline{hamer2025}             & Avaliação de SLAM         & Parcial  & $\times$ & $\times$ & \checkmark & $\times$ \\
\citeonline{maset2022}       & Repetibilidade 3D         & \checkmark & $\times$ & $\times$ & \checkmark & $\times$ \\
\citeonline{amigoni2010}           & Metodologia experimental  & $\times$ & $\times$ & $\times$ & Parcial  & N/A \\
\citeonline{guevaravega2024}       & Artefatos de pesquisa (ES) & $\times$ & N/A & $\times$ & N/A & N/A \\
\textbf{Esta dissertação}     & Framework experimental    & \checkmark & \checkmark & \checkmark & \checkmark & \checkmark \\
\bottomrule
\end{tabularx}
{\footnotesize \checkmark = presente; $\times$ = ausente; Parcial = presente com limitações}
\end{table}

\section{Lacuna Identificada}
\label{sec:rel_lacuna}

A revisão da literatura conduzida nesta dissertação (Seções~\ref{sec:rel_nav}
a~\ref{sec:rel_reprodutibilidade}) não identificou, nos trabalhos analisados, a
combinação específica de características do framework proposto --- protocolo
formal de record--replay, suporte a papéis experimentais em frota, coleta
síncrona de sensores heterogêneos com QoS diferenciada, análise RMSE pairwise
automatizada e interface web sem ROSBridge. Essa afirmação reflete o escopo da
revisão narrativa/comparativa realizada neste capítulo, não uma revisão
sistemática com protocolo de busca formal; não é, portanto, uma alegação de
ineditismo absoluto, mas um relato preciso do que a bibliografia examinada
cobre ou deixa de cobrir.

Essa lacuna não é acidental. Ela reflete uma escolha implícita da comunidade de
robótica experimental: priorizar a contribuição algorítmica (um novo controlador,
um novo algoritmo de SLAM) sobre a infraestrutura de avaliação. Esta dissertação
inverte essa prioridade: a infraestrutura de avaliação \textit{é} a contribuição.

\section{Considerações Finais}
\label{sec:rel_conclusao}

A revisão confirma que os componentes técnicos necessários existem e são maduros
--- Nav2, SLAM Toolbox, bags ROS~2 --- mas não estão integrados de forma coerente
em um framework orientado a experimentos reprodutíveis. O próximo capítulo apresenta
a arquitetura que realiza essa integração, com decisões de design motivadas pelos
problemas identificados na literatura.
